\documentclass{article}
\usepackage[utf8]{inputenc}
\usepackage{listings}
\usepackage{color}

\definecolor{codegray}{rgb}{0.95,0.95,0.95}
\lstset{
  backgroundcolor=\color{codegray},
  basicstyle=\ttfamily\footnotesize,
  frame=single,
  language=Python,
  keywordstyle=\color{blue},
  commentstyle=\color{gray},
  showstringspaces=false
}

\title{Boblesortering i Python}
\author{}
\date{}

\begin{document}
Halla Silje!\\
For likevektskonstannten er det  sånn at når det er likevekt så er:\\
$$K_2=\frac{[H^+]\cdot [CO_3^{2-}]}{[HCO_3^-]}$$
Ved likevekt er er venstre og høyr3e side like. Vi "deler" på $[CO_3^{-2}]$ på begge sider. Da får vi utrykket vi skal se på:\\
$$\frac{K_2}{[CO_3^{2-}]}=\frac{H^+}{[HCO_3^-]}$$
Er det likevekt så vil fortsatt venstre og høyre side være lik, for vi har jo bare snudd på ligningen.\\\\
\textbf{Når $[CO_3^{2-}]>>K_2$:}\\\\
Da blir venstre side av ligningen mindre enn høyre.  dette kan motvirkes ved å redusere $[CO_3^{2-}]$.  likevekten forksyves altså mot venstre, $[CO_3^{2-}]$ avtar.\\\
\textbf{Når $[CO_3^{2-}]<<K_2$:}\\\\
Da blir venstre side av ligningen større enn høyre.  dette kan motvirkes ved å øke $[CO_3^{2-}]$.  likevekten forksyves altså mot høyre, $[CO_3^{2-}]$øker.

\end{document}